\documentclass[11pt,a4paper]{article}
% Preamble for the smc2fc framework documentation.
\usepackage[utf8]{inputenc}
\usepackage[T1]{fontenc}
\usepackage{lmodern}
\usepackage{amsmath, amssymb, amsthm}
\usepackage{algorithm}
\usepackage{algpseudocode}
\usepackage{booktabs}
\usepackage{longtable}
\usepackage{graphicx}
\usepackage{hyperref}
\usepackage{url}
\usepackage{geometry}
\usepackage{xcolor}
\usepackage{listings}
\usepackage{bm}
\usepackage{tikz}
\usetikzlibrary{positioning, arrows.meta, shapes.geometric, fit}

\geometry{margin=1in}

\hypersetup{
    colorlinks=true,
    linkcolor=blue,
    filecolor=magenta,
    urlcolor=cyan,
}

\theoremstyle{definition}
\newtheorem{definition}{Definition}[section]
\newtheorem{proposition}{Proposition}[section]
\newtheorem{theorem}{Theorem}[section]
\newtheorem{remark}{Remark}[section]

% Notation matching old doc so cross-reference works.
\newcommand{\thetadyn}{\bm{\theta}_{\rm dyn}}
\newcommand{\thetactrl}{\bm{\theta}_{\rm ctrl}}
\newcommand{\Lhat}{\widehat{L}}
\newcommand{\xhat}{\widehat{x}}
\newcommand{\dd}{\mathrm{d}}
\newcommand{\E}{\mathbb{E}}
\newcommand{\R}{\mathbb{R}}
\newcommand{\Pcal}{\mathbb{P}}
\newcommand{\Ucal}{\mathcal{U}}
\newcommand{\Xcal}{\mathcal{X}}
\newcommand{\Ncal}{\mathcal{N}}
\newcommand{\Phimax}{\Phi_{\rm max}}
\newcommand{\Fmax}{F_{\rm max}}
\newcommand{\sigmoid}{\mathrm{sigmoid}}
\DeclareMathOperator*{\argmin}{argmin}
\DeclareMathOperator*{\argmax}{argmax}

% Code listing style
\definecolor{codegreen}{rgb}{0,0.6,0}
\definecolor{codegray}{rgb}{0.5,0.5,0.5}
\definecolor{codepurple}{rgb}{0.58,0,0.82}
\definecolor{backcolour}{rgb}{0.97,0.97,0.95}
\lstdefinestyle{code}{
    backgroundcolor=\color{backcolour},
    commentstyle=\color{codegreen},
    keywordstyle=\color{magenta},
    numberstyle=\tiny\color{codegray},
    stringstyle=\color{codepurple},
    basicstyle=\ttfamily\footnotesize,
    breakatwhitespace=false,
    breaklines=true,
    captionpos=b,
    keepspaces=true,
    numbers=left,
    numbersep=5pt,
    showspaces=false,
    showstringspaces=false,
    showtabs=false,
    tabsize=2,
    language=Python,
}
\lstset{style=code}


\title{A Joint LEAN4-and-Julia Charter\\
  for the FSA + \texttt{smc2fc} Stack\\
  \large Part I: LEAN4 as the formal source of truth for FSA models.\\
  Part II: Julia as the production framework for inference and control.\\
  Together: a typed pipeline from LaTeX equations to running code,\\
  with no Python pattern-matching surface for LLM agents to fall through.}
\author{Ajay Talati \quad+\quad Claude (post-mortem author)}
\date{\today}

\begin{document}
\maketitle

\begin{abstract}
This document is a two-part post-mortem-and-charter for a project
that has been burned once too often by LLM coding agents pattern-
matching their way through Python+JAX scientific code. \textbf{Part I}
records the failure mode that made the FSA-v5 model implementation
and its import into the \texttt{smc2fc} framework a multi-day,
expensive process: \emph{LLM coding agents pattern-match on variable
names instead of building a semantic model of types, roles,
dimensions, and pipeline position}. It prescribes the structural
fix on the model side: LEAN4 as the formal source of truth, with
Python as a secondary, differentially-tested artefact.
\textbf{Part II} extends the same diagnosis to the
\texttt{smc2fc} framework itself --- also Python, also
agent-touched, also exposed to the same pattern-matching surface
--- and prescribes the matching structural fix: a Julia port
optimised for Julia idioms (multiple dispatch, parametric types,
native loops, \texttt{AbstractArray} dispatch), with both CPU and
GPU backends in scope from day one. Part II is explicitly \emph{not}
a Python-in-Julia re-skin; an idiomatic Julia codebase is the
deliverable, and the work amounts to roughly seven hours of focused
translation between two scientific computing languages. The two
parts together form a single typed pipeline: LaTeX equations →
LEAN4 reference (model side) → Julia framework (\texttt{smc2fc}
side), with Python remaining in service as the differentially-tested
fast path. This document exists so that no future iteration
relitigates the question.
\end{abstract}

\tableofcontents

\newpage
\part{LEAN4-First Charter for FSA Model Development}
\label{part:lean4}

\section*{Part I scope note}

What follows in Part I is the standalone LEAN4 charter, included
here verbatim so this document is self-contained for an agent who
has not read the original at
\path{FSA_model_dev/LaTex_docs/lean4_first_charter.pdf}. Part I
addresses the \textbf{model side}: how new mathematical models
(FSA-v5 today, future variants tomorrow) are specified and
verified before any Python implementation is written. Part II
(below) extends the same discipline to the \texttt{smc2fc}
framework that consumes those models.

\section{The problem in one paragraph}
\label{sec:problem-lean4}

The current FSA development cycle is: write the mathematical model
in LaTeX; ask an agent to port it to Python; spend days hunting bugs
the agent introduced; eventually ship, often with errata documented
post-deployment. The bottleneck is \emph{agent pattern-matching}, not
modelling complexity. A junior human collaborator with the same LaTeX
source would not make the same class of mistake, because they would
read the equations, internalise that $\theta_{\rm ctrl} \in
\mathbb{R}^{16}$ and $\theta_{\rm dyn} \in \mathbb{R}^{37}$ are
different objects, and refuse to compare them as scalar intervals.
LLM agents reach for the strongest training association of the name
\texttt{theta} and write the most-likely-looking code. The fix has
to be structural --- something that makes the wrong code refuse to
compile --- not motivational.

The audience for this document is two-fold: (a) the user, as a
personal accountability artefact recording why the LEAN4-first
process is being adopted and at what point; and (b) every future
Claude (or other LLM) coding agent that touches an FSA model, who
must read this charter at session start and follow the process it
prescribes.

\section{Concrete evidence}
\label{sec:evidence}

This section anchors the post-mortem in named, citable incidents.
No vague complaints; every claim has a file path or commit hash
attached.

\subsection{Variable-name collisions in the live code}
\label{sec:evidence-collisions}

The same identifier \texttt{theta} refers to two fundamentally
different mathematical objects in two parts of the codebase, and
each carries a different dimension, role, and domain.

\paragraph{(a) \texttt{theta} as 16-D RBF schedule coefficients
(controller decision variable).} In
\path{models/fsa_high_res/control.py}, lines 61--65:
\begin{lstlisting}[language=Python]
def schedule_from_theta(theta: jnp.ndarray) -> jnp.ndarray:
    theta_B = theta[:n_anchors]
    theta_S = theta[n_anchors:]
    raw_B = c_Phi + jnp.einsum('a,ta->t', theta_B, Phi_design)
\end{lstlisting}
Here \texttt{theta} is a length-16 vector ($2 \times n_{\rm anchors}$,
with $n_{\rm anchors} = 8$), interpreted as RBF basis coefficients
that, after a sigmoid envelope, parameterise the dual-channel
control schedule $\Phi(t)$. It is the controller's
\emph{decision variable}. It lives in $\mathbb{R}^{16}$. Its prior
breadth is the scalar \texttt{sigma\_prior} in
\path{smc2fc/control/config.py:28}.

\paragraph{(b) \texttt{theta} as 37-D dynamics-parameter posterior
(SMC$^2$ inference target).} In
\path{models/fsa_high_res/estimation.py}, the function
\verb|get_init_theta()| at line 513--520:
\begin{lstlisting}[language=Python]
def get_init_theta():
    """Return prior-mode parameter vector (length = len(PARAM_PRIOR_CONFIG))."""
    return np.array([_mode(pt, pa) for _, (pt, pa)
                     in PARAM_PRIOR_CONFIG.items()],
                    dtype=np.float32)
\end{lstlisting}
Here \texttt{theta} is a length-37 vector of physiological
parameters --- rate constants, amplitudes, decay times, observation
noise scales --- whose posterior the SMC$^2$ machinery infers from
data. It is the inference \emph{target}, not a decision variable.
It lives in $\mathbb{R}^{37}$. Its prior is component-wise
log-normal, not Gaussian; the prior breadth varies per parameter.

\paragraph{Why this matters.} An agent reading either file sees
\texttt{theta}, a flat numpy array, and reaches for ``Gaussian prior
on $\theta$'' --- which is correct in case (a) and wrong in case
(b). An agent reasoning across files conflates the two and treats
``the prior on $\theta$'' as a single concept with a single dimension
and a single breadth. There is no warning sign in the Python code
itself; the type is just \texttt{jnp.ndarray} in both places.

\subsection{The Bug 5 incident}
\label{sec:bug5}

The named incident that motivated this charter, recorded verbatim
from the user's feedback memory at
\path{~/.claude/projects/.../memory/feedback_no_fabricated_bugs.md}:

\begin{quote}\itshape
During the FSA-v5 verification session on 2026-05-06 I logged a
``Bug 5'' that claimed $\theta \in [-0.4, +0.5]$ for the controller's
RBF schedule prior to keep the system in the healthy island.
$\theta$ is a 16-dim vector (RBF coefficients across both Phi
channels and \texttt{n\_anchors=8} anchors). The healthy island is a
2D region in $(\Phi_B, \Phi_S)$ \textbf{control output space}, not in
$\theta$ space. Comparing them as scalar intervals is incoherent. I
conflated:
\begin{itemize}
\item SMC$^2$ parameter posterior $\theta$ (37-dim, the inferred
  parameters)
\item Controller RBF schedule $\theta$ (16-dim, the decision
  variables)
\end{itemize}
both of which use \texttt{sigma\_prior} in the bench code, mentally
collapsed to a single scalar. The bench code in front of me had
\texttt{theta\_B = theta[:n\_anchors]} --- explicit vector indexing
--- and I didn't internalize it.
\end{quote}

The agent in question was me. The user's verdict was ``stupid
mistake''; the request was for guardrails. This document is the
guardrail.

\subsection{Plan churn in \texttt{claude\_plans/}}
\label{sec:plan-churn}

The repository's \path{claude_plans/} directory accumulates one
markdown file per plan-mode session. The CLAUDE.md convention is
that meaningful updates to a plan add a \verb|> Updated:| audit
line below the title. As of this writing the directory contains
eight plan files, of which \emph{seven have zero \texttt{Updated:}
lines} --- they were drafted, never executed, never closed:

\begin{itemize}\small
  \item \path{FSA_v5_Implementation_Plan_for_the_fsa_high_res_Folder_2026-05-04_2238.md} (0 updates).
  \item \path{Lyapunov_Stability_Analysis_for_FSA_v4_Plan_2026-05-04_1448.md} (0 updates).
  \item \path{Ajays Notes Plan to update to version 5.md} (0 updates).
  \item \path{Advice_on_smc2fc_Port_Cost_Function_Architecture_2026-05-05_1055.md} (0 updates --- advice-only, never tied to a closed task).
  \item \path{Techical documentation plan - FSA_version_5_technical_guide tex.md} (0 updates).
  \item \path{Ajays Notes - 1.md}, \path{Ajays Notes - 2.md} (0 updates).
\end{itemize}

The single plan that did execute is
\path{FSA_dev_sandbox_mirror_the_SWAT_model_dev_pattern_v4_branch_2026-05-04_2353.md}
(1 update line, 13/13 tests green at completion). The ratio of
abandoned to executed plans (7:1) is not a sign of laziness; it is
a sign that planning sessions repeatedly re-derive the same
mathematical model in slightly-different shapes because no formal
artefact pins the model down.

\subsection{Errata after deploy}
\label{sec:errata}

Commit \texttt{d886f36} on the v4 branch:
\begin{quote}\itshape
docs: add FSA-v4 errata, open-loop limitation, and SMC$^2$ vs OT
comparison.

Documents the gap between the §4 narrative and the actual behaviour
of \texttt{tools/optimize\_dynamic\_plan\_v4.py}, then frames smc2fc
as the structural remedy.
\end{quote}
A documentation-correction commit recording bugs caught
\emph{after} deployment, not during development. The fact that
errata was needed at all is the evidence that the in-development
verification was not strong enough; it does not matter how many
tests pass on the final commit if the bugs are documented in a
follow-up commit instead of caught at the place they were
introduced.

\section{Root-cause analysis}
\label{sec:rca-lean4}

\paragraph{LLMs are surface pattern-matchers.} They read source
code primarily as text, not as a typed graph of objects with
roles. When an agent sees \texttt{theta}, the most-trained
association fires; the surrounding context is treated as
secondary. This is not a defect of any particular model; it is the
structural consequence of how the system is built. Asking an agent
to ``be more careful'' is asking it to override its own training
prior every time --- which empirically does not happen reliably,
as §\ref{sec:bug5} demonstrates.

\paragraph{Python provides no formal scaffolding to compensate.}
\texttt{jnp.ndarray} is the type of both the 16-D control
coefficient vector and the 37-D parameter posterior. There is no
shape-aware static analyser running in the loop, no role tagging
in the type system, no proof obligation attached to a function
signature that the inputs satisfy any pre-condition. The senior
engineer (the user) ends up as the only line of defence, and the
\pounds\pounds\pounds\ in tokens is the cost of the user catching agent errors after
the fact.

\paragraph{Why ``tell agents to be more careful'' doesn't work.}
The user's feedback memory at \path{feedback_no_fabricated_bugs.md}
encodes a PRIMARY RULE --- \emph{build a semantic model BEFORE any
claim} --- which is correct as a principle but unenforceable as a
process. The rule depends on the agent \emph{choosing} to think
rather than pattern-match. §\ref{sec:bug5} is the empirical proof
that the choice is not reliably made. Every agent session is one
roll of the dice on whether the rule fires that day.

\section{The structural fix: LEAN4 as the source of truth}
\label{sec:fix}

\subsection{Dependent types make the collision a compile error}

A LEAN4 function with the signature
\[
\texttt{evaluateCost} : \texttt{RBFCoeffs}\;16 \to \mathbb{R}
\]
syntactically refuses any value not of the exact type
\texttt{RBFCoeffs 16}. The 37-dimensional parameter posterior
$\theta_{\rm dyn} : \texttt{ParameterVector}\;37$ is a different
type, and \texttt{lake build} will reject any code that conflates
them. Pattern-matching on names becomes harmless because the type
system catches the consequence at compile time, before any
incoherent reasoning reaches the user.

The two-dimensional output region $(\Phi_B, \Phi_S)$ from
§\ref{sec:bug5} would, in LEAN4, have a type like
\texttt{ControlOutput} that is \emph{not} the type of any prior on
$\theta$. Writing $\theta \in [-0.4, +0.5]$ would not compile, full
stop, because the right-hand side is a scalar interval and the
left-hand side is a vector. The agent could not have made the
mistake.

\subsection{LEAN4 is also a programming language}

LEAN4 compiles to native code via C. \texttt{lake build} produces
a binary; a thin Python \texttt{ctypes} wrapper exposes each
verified function as a Python callable. The same LEAN4 file is
both the formal specification and the executable reference
implementation. There is no ``the spec drifted from the code''
failure mode because the spec \emph{is} the code.

\subsection{mathlib already has the infrastructure}

The mathematical objects FSA-v5 needs ---
finite-dimensional real vector spaces, matrix operations,
exponentials and logarithms, Hill functions, sigmoids,
piecewise-linear schedule decoders, ODE integration --- are all
already library imports in \texttt{mathlib}. There is no novel
mathematics to formalise. Transcribing the FSA-v5 LaTeX into LEAN4
is genuinely just translation.

\subsection{What this does not solve}

\begin{itemize}
\item \emph{Floating-point semantics.} The gap between the LEAN4
  notion of $\mathbb{R}$ and IEEE-754 double-precision arithmetic
  is real and not addressed by this charter; it is bounded by the
  tolerance threshold in the differential test (§\ref{sec:process}
  step 6).
\item \emph{GPU correctness.} The differential test runs the
  Python implementation on whatever backend it normally uses
  (JAX/XLA, GPU); the LEAN4 binary is CPU-only. Disagreements
  driven by GPU-specific behaviour are out of scope for the LEAN4
  reference itself and trusted as before.
\item \emph{The transcription itself.} If the LaTeX$\to$LEAN4
  transcription is wrong, both the LEAN4 reference and (likely)
  the agent-written Python will be wrong in the same way. This is
  why the LaTeX$\to$LEAN4 step is human-reviewed line-by-line; it
  is the one place reviewer attention has to land.
\end{itemize}

\section{The new development process}
\label{sec:process}

The standing rule going forward, and the charter every future
agent must follow:

\begin{quote}\bfseries
For every new mathematical model, every model extension, and every
import into \texttt{smc2fc}, the LEAN4 specification is written
before any Python implementation. The LEAN4 native binary is the
authoritative reference. Python is differentially tested against
it. Disagreement beyond a tolerance threshold is by construction a
Python bug.
\end{quote}

Per function, the workflow is:

\begin{enumerate}
\item \textbf{LaTeX equations (as today).} Drift, diffusion, cost,
  plant, schedule decoder, observation model --- written in
  \path{LaTex_docs/sections/} using the existing notation
  conventions.
\item \textbf{LEAN4 transcription.} A new \path{lean/} directory in
  the repo. The LaTeX equations are transcribed into LEAN4
  \emph{definitions} (typed, with dimensions encoded in the type)
  and \emph{implementations} (runnable code). Line-by-line,
  side-by-side with the LaTeX. Straight translation, no design
  work, no novel mathematics.
\item \textbf{LEAN4 native compilation.} \texttt{lake build}
  produces a native binary. A Python \texttt{ctypes} wrapper in
  \path{lean/python_bridge/} exposes each function as a Python
  callable.
\item \textbf{Python implementation (agent-written, fast path).}
  The agent writes the JAX/Python implementation as today,
  optimised for GPU. The agent is permitted to pattern-match;
  step 6 catches the consequences.
\item \textbf{(Optional) property proofs.} Lipschitz, basin-of-
  attraction theorems, monotonicity --- the kind of statement that
  makes the LEAN4 spec a \emph{proof} not just a typed transcription.
  Defer unless a specific incident motivates one. The charter does
  not gate on these.
\item \textbf{Differential testing.} A \texttt{pytest} harness using
  \texttt{hypothesis} for random-input generation evaluates BOTH the
  LEAN4 binary AND the Python implementation on a battery of inputs
  sampled from the physiological-bound box. Any disagreement beyond
  a tolerance threshold (default $10^{-6}$ in single-step drift,
  $10^{-4}$ in integrated trajectory) is a Python bug.
\item \textbf{\texttt{smc2fc} port.} The port is a thin shim from
  the verified LEAN4 reference into the framework's
  \texttt{EstimationModel} / \texttt{SDEModel} interfaces. The
  shim is also differentially tested against the LEAN4 binary
  before being merged.
\end{enumerate}

\section{Scope and timing (model side)}
\label{sec:scope-lean4}

\textbf{This is translation work, not research.} The LaTeX spec
already exists; LEAN4 has the linear-algebra and ODE infrastructure;
\texttt{mathlib} supplies the lemmas. The full FSA-v5 model is
\textbf{one day's work for a competent human developer}. Any
estimate quoting weeks per function is inflation; the user has
explicitly rejected that framing and this document is bound by the
rejection.

\begin{center}\small
\begin{tabular}{lll}
\toprule
Function & File & LEAN4 transcription \\
\midrule
FSA-v5 drift              & \path{models/fsa_high_res/_dynamics.py} & $\sim 2$ hours \\
Cost (chance-constrained) & \path{models/fsa_high_res/control_v5.py} & $\sim 1.5$ hours \\
StepwisePlant             & \path{models/fsa_high_res/_plant.py} & $\sim 1$ hour \\
Schedule decoder (RBF)    & \path{models/fsa_high_res/control.py} & $\sim 30$ min \\
Observation model         & (in \path{estimation.py}) & $\sim 1$ hour \\
FFI wrapper + diff tests  & \path{lean/python_bridge/} & $\sim 2$ hours \\
\midrule
Total                     & --- & $\sim 8$ hours \\
\bottomrule
\end{tabular}
\end{center}

The LEAN4 source is roughly the same length as the LaTeX it
transcribes. Most of the time goes into wiring up the FFI bridge
and writing the differential-test harness; the math itself is
direct.

The cost/value frame: a single 8-hour LEAN4 transcription pass
amortises across \emph{every future FSA version, every future model
extension, and every future agent attempting a port to
\texttt{smc2fc}}. The alternative --- paying for repeated agent
rounds of pattern-match-then-hunt-bugs --- is what burned the
\pounds\pounds\pounds\ this charter exists to stop.

\section{What is and isn't covered (model side)}
\label{sec:in-out-lean4}

\paragraph{In scope.} Drift, cost, plant, schedule decoder,
observation model. Typed LEAN4 reference implementations plus
differential tests. The agent-written FSA-v5 surface in
\path{models/fsa_high_res/} (2{,}200 Python LOC across 4 files;
the LEAN4 transcription is much shorter).

\paragraph{Out of scope (still trusted).} The SMC$^2$/PF framework
code in \texttt{smc2fc/core/} and \texttt{smc2fc/filtering/} is
already tested upstream and is not subject to \emph{this Part}.
\textbf{Note:} Part II below proposes a Julia port of the framework
that does subject this code to type-safety scrutiny by structural
means; that is the joint deliverable of this combined charter.
HMC and NUTS internals, the JAX-native compile-once architecture,
GPU correctness, and IEEE-754 floating-point semantics are
trusted as before.

\paragraph{Trusted by transcription.} The LaTeX$\to$LEAN4 step is
human-reviewed line-by-line. This is the one place reviewer
attention has to land; the rest of the process is mechanical.

\section{Audit checklist (model-side agent)}
\label{sec:checklist-lean4}

A future agent runs this checklist at session start, before
touching any FSA-related Python file:

\begin{itemize}
\item[$\square$] Is there a LEAN4 reference for the function I am
  about to touch? \emph{If no:} the charter says I am not allowed
  to write the Python until the LEAN4 reference exists. Stop and
  flag.
\item[$\square$] Does \texttt{lake build} succeed on the current
  commit (the LEAN4 reference type-checks)?
\item[$\square$] Does the Python differential test pass against
  the LEAN4 binary on the current commit (\texttt{pytest
  tests/test\_lean4\_diff.py} is green)?
\item[$\square$] Have I written down --- in plain English ---
  the type, role, and pipeline position of every variable I am
  about to reason about? (Per the PRIMARY RULE in
  \path{feedback_no_fabricated_bugs.md}.)
\item[$\square$] If I am about to claim a bug: do I have a
  reproducer (input $\to$ both implementations $\to$ disagreement
  beyond the tolerance threshold)? \emph{If no:} it is not a bug,
  it is an open question; record it as such and move on.
\end{itemize}

If any box is unchecked, the agent does not proceed. The session
is paused; the user is notified; the missing artefact is produced
before any code is changed.

\section{Closing of Part I}
\label{sec:closing-lean4}

The point of Part I is accountability on the model side. The named
incident in §\ref{sec:bug5} was caused by an agent (me) pattern-
matching on \texttt{theta} across two architectural layers it does
not unify. The cost was burned trust and burned time. The charter
prescribed in §\ref{sec:process} makes that class of error
structurally impossible: the LEAN4 type-checker is unforgiving, and
disagreements between LEAN4 and Python are by construction Python
bugs.

LEAN4 is never wrong about types. Python is. Agents pattern-match;
LEAN4 does not. Part I is the bridge that makes the agent's
limitations harmless on the model side.

\textbf{Bridge to Part II.} The LEAN4 charter solves half of the
problem --- the model side. What about the framework
\texttt{smc2fc} itself? It is also Python, also agent-touched,
also exposed to the same pattern-matching surface, and the
empirical evidence in §\ref{sec:evidence} (the \texttt{theta} 16-D
vs 37-D collision, the closure-of-data recompile bug, the
NUTS-inside-vmap-inside-\texttt{lax.while\_loop} cliff) is at least
as pointed about \texttt{smc2fc} internals as it is about
\texttt{models/fsa\_high\_res/}. Part II prescribes the matching
structural fix for the framework side.

\newpage
\part{Julia Port Charter for the \texttt{smc2fc} Framework}
\label{part:julia}

\section{The problem in one paragraph (framework side)}
\label{sec:problem-julia}

The \texttt{smc2fc} codebase is doing the right mathematics; the
failure mode is the same as the FSA-v5 model side: agent
pattern-matching across a Python codebase with no type-level
enforcement of dimensions or roles. Part I solves half of this ---
the model side --- by making LEAN4 the formal source of truth for
the math \texttt{smc2fc} consumes. Part II solves the other half
by porting the framework itself to a language where the type system
catches the pattern-matching errors at compile time. The audience
is the same two populations: the user (accountability) and future
LLM agents (constitution loaded at session start).

\section{The Python + LLM perfect storm}
\label{sec:storm}

Python and LLM coding agents combine into a particularly bad
failure surface for the kind of math \texttt{smc2fc} does. This
section names the four failure modes inside \texttt{smc2fc/}
specifically.

\subsection{Everything is \texttt{jnp.ndarray}}
\label{sec:storm-ndarray}

Python+JAX has one type for every numerical object regardless of
shape, role, or origin. Concrete signatures from the live code:

\paragraph{(a)} \path{smc2fc/core/jax_native_smc.py:67} ---
\begin{lstlisting}[language=Python]
def _ess_at_delta(loglikelihood_fn, particles, delta):
    ...
\end{lstlisting}
\texttt{particles} is a 2-D array of shape \texttt{(n\_smc, d)}; the
type annotation is just \texttt{jnp.ndarray}. An agent passing a
single particle (a 1-D array of shape \texttt{(d,)}) gets a
runtime broadcast surprise --- the function returns a wrong-shape
result silently.

\paragraph{(b)} \path{smc2fc/filtering/_gk_kernel.py:185} ---
\begin{lstlisting}[language=Python]
def smooth_resample_basic(particles, log_weights, ...):
    ...
\end{lstlisting}
\texttt{log\_weights} could be a 1-D vector of length \texttt{n\_pf}
(if the caller is operating on a single particle filter) or a 2-D
\texttt{(n\_smc, n\_pf)} block (if the caller has already
\texttt{vmap}'d over the SMC$^2$ outer cloud). Nothing in the
signature says which. The function happens to work in both cases
because of broadcasting; an agent reasoning about either case
internalises one of the two and writes downstream code that breaks
in the other.

\paragraph{(c)} \path{smc2fc/core/tempered_smc.py} --- the
BlackJAX-wrapped backend captures \texttt{data} in a closure per
window. This is the bug §6 of the framework documentation
documents: per-window recompiles dominate the runtime. The Python+
JAX foible is invisible at the call site; an agent reading the
function signature has no way to predict that \texttt{data} ends up
hashed into the JIT cache key.

\subsection{Mutability and pure-functional code coexist invisibly}

In \path{smc2fc/control/diagnostics.py} the \texttt{matplotlib}
helpers are imperative side-effectful Python; in
\path{smc2fc/core/jax_native_smc.py} the SMC loop is
pure-functional for JIT compatibility. Both files use the same
\texttt{def} keyword, the same \texttt{jnp.ndarray} types, the same
import style. An agent editing one file reaches for the pattern of
the other and silently breaks the JIT contract --- mutating an
array inside what was a pure function, capturing a non-traceable
Python object inside a JIT'd closure. The LEAN4 charter's PRIMARY
RULE (``build a semantic model before any claim'') has to fire on
every cross-file edit; empirically it does not.

\subsection{Closures over data are a recompile landmine}

A long-running session that thinks it is reusing a JIT-compiled
function silently re-traces because a dataclass's \texttt{\_\_hash\_\_}
changed. Documented in framework-doc §6.1; an agent cannot infer
this from any signature in the file because the cause lives in
JAX's tracing semantics, not in the Python source.

\subsection{The vmap/while-loop nesting cliff}
\label{sec:storm-vmap}

§6.4 of the framework documentation records that NUTS-inside-
\texttt{vmap}-inside-\texttt{lax.while\_loop} is 60--85$\times$
slower than HMC because of dynamic-tree-building inside nested
control flow. This is a Python+JAX architectural consequence, not a
mathematical one --- and it is exactly the kind of trap an agent
walks into by pattern-matching ``NUTS is a strict upgrade from HMC''
from training data. The signature
\verb|nuts.build_kernel(...)| says nothing about the nesting cliff;
the agent has to know JAX's compilation model to predict the
disaster.

\section{Why Julia narrows the failure surface}
\label{sec:why-julia}

\subsection{Multiple dispatch + parametric types}

A Julia function declared
\begin{lstlisting}
function ess_at_delta(loglik::F, particles::Matrix{Float64},
                      delta::Float64) where {F<:Function}
    ...
end
\end{lstlisting}
refuses anything else at compile-time. \texttt{Vector\{Float64\}}
and \texttt{Matrix\{Float64\}} are different types; passing one
where the other is expected does not silently broadcast. The
single-particle vs particle-cloud confusion from
§\ref{sec:storm-ndarray}(a) becomes a method-not-found error at the
call site, not a wrong-shape result downstream.

\subsection{Functional style is the default}

Mutation is opt-in via the \texttt{!} suffix convention:
\texttt{resample!(particles, weights)} mutates in place;
\texttt{resample(particles, weights)} returns a fresh array. An
agent reading any function signature can tell instantly which
mode it is in. The Python invisibility from §\ref{sec:storm}'s
``imperative diagnostics file vs pure-functional SMC loop file''
disappears.

\subsection{Code reads like the LaTeX}

A Julia drift function with parametric types reads almost
identically to the equation it transcribes:
\begin{lstlisting}
function drift(x::FSAState{6}, theta::DynParams)::FSATangent{6}
    ...
end
\end{lstlisting}
This is much closer to the LEAN4 spec from Part I than the
equivalent Python. An agent moving between LEAN4 and Julia is
moving between two flavours of typed dispatch; an agent moving
between LEAN4 and Python is moving between a typed system and an
untyped one and has to translate types out of sight at every step.

\subsection{No two-language problem}

Python+JAX exists because Python is too slow for the inner loop;
JAX is the workaround that ships the inner loop to XLA. Julia is
fast natively. The JIT-traced / non-traced split that creates the
invisible mutability landmine of §\ref{sec:storm} is a
\emph{Python} artefact; in Julia, the inner and outer loops are
the same kind of code.

\subsection{Static analysis catches the rest}

\texttt{JET.jl} runs a static type analysis on the entire package
and flags every type uncertainty. This is a single command:
\begin{lstlisting}
julia --project=. -e 'using JET; report_package(SMC2FC)'
\end{lstlisting}
There is no analogue in Python (\texttt{mypy} doesn't reach into
JAX trees, and even where it does, the pervasive
\texttt{jnp.ndarray} type defeats it).

\subsection{LEAN4 alignment}

A LEAN4 specification for the FSA model (per Part I) translates
to Julia almost one-to-one because both languages are
dispatch-based and type-driven. The LEAN4 $\to$ Julia
transcription is shorter than LEAN4 $\to$ Python by roughly the
ratio of the language verbosities. The combined LEAN4 reference
plus Julia framework forms a single coherent typed pipeline
\emph{from the LaTeX equations to the running code}. Python sits
awkwardly in the middle of this pipeline today; remove it from
the framework side and the pipeline aligns.

\section{Existing Julia scaffolding --- what we re-use}
\label{sec:scaffolding}

The bulk of \texttt{smc2fc/} is wiring well-known algorithms
together. Julia already has each of those algorithms as a
maintained library; the port is mostly \emph{replacing custom code
with library calls}. The table below is the per-module mapping.

\begin{center}\footnotesize
\begin{longtable}{p{4.5cm}p{4.5cm}p{6cm}}
\toprule
Python \texttt{smc2fc} module & Julia replacement & Notes \\
\midrule
\endhead
\path{core/jax_native_smc.py} HMC kernel & \texttt{AdvancedHMC.jl} & The original HMC implementation BlackJAX modelled on. Direct drop-in. \\
\path{core/tempered_smc.py} & Custom (built on \texttt{AbstractMCMC.jl}) & No exact analogue but the abstractions are there; $\sim 150$ lines of Julia. \\
\path{core/sf_bridge.py} BW geodesic & \texttt{LinearAlgebra} (stdlib) + \texttt{Distributions.jl} & Matrix square root, multivariate Gaussian. $\sim 50$ lines. \\
\path{core/mass_matrix.py} & \texttt{Statistics.var} (stdlib) & Literal one-liner. \\
\path{filtering/sinkhorn.py} + \path{transport_kernel.py} + \path{resample.py:ot_resample_lr} & \texttt{OptimalTransport.jl} & Maintained, supports entropic + low-rank. \\
\path{filtering/_gk_kernel.py} Liu--West & Custom ($\sim 80$ lines) & Standard moment-match + Gaussian smoother; trivial in Julia. \\
\path{filtering/gk_dpf_v3_lite.py} & Custom ($\sim 150$ lines) & Bootstrap filter loop; reuses \texttt{Distributions.jl} for log-density. \\
\path{simulator/sde_solver_diffrax.py} & \texttt{StochasticDiffEq.jl} (part of \texttt{DifferentialEquations.jl}) & Gold standard. Replaces diffrax wholesale. \\
\path{transforms/unconstrained.py} & \texttt{Bijectors.jl} & The Julia ecosystem's standard bijection library. \\
Einsum-style ops (Liu--West kernels, RBF designs) & \texttt{Tullio.jl} or built-in broadcasting & Often broadcasting alone suffices. \\
Autodiff for HMC gradients & \texttt{Enzyme.jl} (preferred) or \texttt{Zygote.jl} & \texttt{Enzyme.jl} is more efficient for the kind of code here. \\
Static analysis & \texttt{JET.jl} & The ``Julia answer to the LLM pattern-match risk''. \\
Testing harness & \texttt{Test.jl} (stdlib) & Standard. \\
GPU arrays & \texttt{CUDA.jl} & NVIDIA backend; \texttt{CuArray} parametric type plays naturally with the generic \texttt{AbstractArray\{T,N\}} dispatch the rest of the package uses. \\
Cross-device kernels & \texttt{KernelAbstractions.jl} & Write a kernel once, runs on CPU (\texttt{CPU()}) or GPU (\texttt{CUDABackend()}). Used where broadcasting alone won't suffice. \\
\midrule
\multicolumn{3}{l}{\emph{Performance-critical scaffolding (added per external review):}} \\
\midrule
Small fixed-size arrays (state, params) & \texttt{StaticArrays.jl} & \textbf{Critical.} \texttt{SVector\{6, Float64\}} is stack-allocated, zero-overhead, strictly sized. The Julia equivalent of XLA unrolling small arrays into registers. Without it, \texttt{State\{6\}} and \texttt{DynParams} would heap-allocate on every step and the inner SDE/SMC loop would regress badly versus JAX. \\
Particle cloud layout & \texttt{StructArrays.jl} & Struct-of-Arrays (SoA) memory layout. A \texttt{Particle} struct stored in a \texttt{StructArray} reads as \texttt{for p in particles} but the underlying memory is contiguous per field; coalesced GPU access. \\
Parameter blocks & \texttt{ComponentArrays.jl} & Flat \texttt{AbstractVector} view of nested \texttt{DynParams} / \texttt{CtrlParams} with dot-access syntax (\texttt{theta.k\_FB}). Plug-and-play glue between \texttt{AdvancedHMC.jl}, \texttt{Enzyme.jl}, and the domain types. \\
Likelihood / posterior interface & \texttt{DensityInterface.jl} & Standard wrapper that makes any log-density function plug-and-play with \texttt{AdvancedHMC.jl} and the wider Julia MCMC ecosystem. \\
Particle-filter primitives & \texttt{AdvancedPS.jl} (optional) & Battle-tested SMC sweeps, particle-Gibbs samplers, robust resampling schemes from the Turing ecosystem. May replace the $\sim 150$-line custom bootstrap-filter loop in \path{src/Filtering/Bootstrap.jl}. \\
CPU thread parallelism & \texttt{Polyester.jl} (or \texttt{Base.Threads}) & Replaces JAX's \texttt{vmap} on the CPU backend. \texttt{Polyester.@batch} has lower overhead than \texttt{Threads.@threads} for tight inner loops. \\
\bottomrule
\end{longtable}
\end{center}

The custom code in the port --- the SMC$^2$ tempering loop, the
bootstrap filter wrapper, Liu--West shrinkage, the
cost-as-likelihood controller loop, the bridge-warm-start glue ---
is roughly \textbf{1{,}000--1{,}200 lines of Julia}. Compare to the
$\sim 5{,}600$ lines of Python the port replaces. The reduction is
mostly the disappearance of JAX-specific boilerplate (closures,
\texttt{Partial} wrappers, \texttt{lax.while\_loop} scaffolding,
manual \texttt{jit} placement).

\section{Hybrid CPU/GPU execution strategy}
\label{sec:hybrid}

The §\ref{sec:storm-vmap} ``vmap/while-loop nesting cliff'' has a
deeper architectural cause that informs how the Julia port should
allocate work to hardware. SMC$^2$ is structurally
\emph{heterogeneous}: its inner loop and its outer loop have
diametrically opposed hardware preferences, and the Python+JAX
implementation could not separate them because JAX's JIT/closure
model effectively forces all inside-trace work onto the same device.
Julia's per-function compilation lifts that constraint, and the port
should exploit it from day one.

\subsection{The two phases of SMC$^2$ have opposite hardware
profiles}

\paragraph{Inner loop (particle filter, SDE simulation).} Massively
parallel, SIMD-friendly, predictable control flow. Advancing
$N \sim 10^4$ SDE trajectories simultaneously is the textbook
GPU workload. Ideal hardware: \textbf{NVIDIA RTX 5090} (or
equivalent --- the dominant cost is memory bandwidth $\times$ CUDA
core count).

\paragraph{Outer loop (parameter rejuvenation: HMC / NUTS).}
Highly branchy, sequential per chain, dynamic control flow ---
particularly NUTS, where the trajectory length is determined by a
U-turn check whose branch pattern varies per particle. GPUs
suffer from \emph{warp divergence} when neighbouring threads take
different branches; this is exactly what causes the §\ref{sec:storm-vmap}
60--85$\times$ slowdown. Ideal hardware: \textbf{Intel Core
Ultra 9 285K} (or equivalent --- a strong-IPC, deep-cache,
branch-predictor-rich CPU).

\subsection{The Julia hybrid pipeline}

The Julia port should run the two phases on the hardware they want
and connect them with a low-bandwidth interface across the PCIe
bus. The data-flow per outer SMC$^2$ iteration:

\begin{enumerate}
\item \textbf{Phase A (GPU).} The current parameter cloud
  $\{\theta^{(m)}\}_{m=1}^M$ lives on the host. For each
  $\theta^{(m)}$, a particle filter with $N$ state particles
  is run \emph{on the GPU} via
  \texttt{CuArray}-resident state, advancing all particles in
  parallel via \texttt{KernelAbstractions.@kernel} or
  \texttt{StochasticDiffEq.jl}'s GPU dispatch. The output is a
  scalar $\widehat{L}_N(\theta^{(m)})$ per outer particle.
\item \textbf{Phase B (PCIe, summary statistics only).} Only the
  scalar $\widehat{L}_N(\theta^{(m)})$ (and, if HMC, its gradient
  with respect to the 37-D $\theta^{(m)}$) is transferred from
  device to host. Transferring a scalar or a 37-D vector across
  PCIe Gen 5 takes microseconds --- entirely negligible against the
  SDE integration time. The $N \times d$ particle state matrix is
  \emph{not} transferred; it stays GPU-resident.
\item \textbf{Phase C (CPU).} Tempered SMC$^2$ outer loop runs on
  the host. Adaptive $\delta\lambda$ via bisection. Mass-matrix
  re-estimation. HMC (or, optionally, NUTS) rejuvenation moves
  via \texttt{AdvancedHMC.jl} on the CPU, where the dynamic
  control flow benefits from branch prediction. The
  multi-core architecture is exploited by wrapping the outer
  particle loop in \texttt{Threads.@threads} or
  \texttt{Polyester.@batch} --- each thread driving its own
  GPU dispatch for inner work.
\item \textbf{Phase D (PCIe, parameter vectors only).} The new
  parameter proposals $\{\theta'^{(m)}\}$ are sent back across
  PCIe to the GPU. 37-D vectors $\times M$ outer particles is
  again negligible bandwidth.
\item \textbf{Repeat.}
\end{enumerate}

\subsection{Why this was infeasible in Python+JAX}

The JAX backend bakes the log-likelihood into a single JIT-compiled
closure. Inside that closure, any device split has to be expressed
as JAX operations, which means the entire trace must lower to one
backend (XLA-CPU or XLA-GPU). Forcing the inner loop to GPU and
the outer loop to CPU within the same JIT trace would require
host-callback boundaries that defeat the JIT, or a two-program
architecture with high-frequency device handoffs that defeat
XLA's fusion. Either way, the per-iteration cost of the boundary
exceeds the saving from the right-hardware allocation. So the
JAX implementation chose ``all on GPU'', NUTS-inside-vmap-inside-
\texttt{lax.while\_loop}, and paid the §\ref{sec:storm-vmap} 60--
85$\times$ tax.

In Julia, each function compiles independently. The
\texttt{run\_particle\_filter!(filter\_state::GPUFilterState,
theta::CuArray)} method dispatches to a CUDA kernel; the
\texttt{rejuvenate!(cloud::CPUParameterCloud, log\_lik::Vector)}
method dispatches to ordinary CPU code calling
\texttt{AdvancedHMC.jl}. The boundary between them is just two
function calls and one \texttt{Array(scalar\_loglik)} transfer ---
no JIT trace to keep happy, no monolithic compilation unit to
fragment.

\subsection{Architectural type boundaries}

Two parametric types make the hybrid contract explicit:
\begin{lstlisting}
# Particle states: GPU-resident.
struct GPUFilterState{T}
    particles::CuArray{T, 2}      # (n_pf, d_state)
    weights::CuArray{T, 1}        # (n_pf,)
end

# Outer parameter cloud: CPU-resident.
struct CPUParameterCloud{T}
    theta::Matrix{T}              # (n_smc, d_theta), CPU
    log_lik::Vector{T}            # (n_smc,), CPU
end
\end{lstlisting}
The likelihood function consumed by \texttt{AdvancedHMC.jl}
handles the host-device transfer internally:
\begin{lstlisting}
function compute_likelihood(theta_cpu::AbstractVector{Float64},
                            filter_state::GPUFilterState)
    theta_gpu = CuArray(theta_cpu)              # 37-D, microseconds
    log_lik_gpu = run_particle_filter!(filter_state, theta_gpu)
    return Array(log_lik_gpu)[1]                # scalar back, microseconds
end
\end{lstlisting}
When \texttt{Enzyme.jl} or \texttt{Zygote.jl} differentiates this
for HMC gradients, the VJP also handles the host-device boundary
correctly --- the autodiff stays inside whichever device is doing
the work and only the gradient w.r.t.\ \texttt{theta\_cpu}
crosses back.

\subsection{What this does to the NUTS question}

§\ref{sec:storm-vmap} records that NUTS-inside-vmap-inside-
\texttt{lax.while\_loop} is the architecturally infeasible
combination, not NUTS itself. Under the hybrid pipeline, NUTS runs
on the CPU where its branchy dynamic-tree-building is well-
matched to the hardware. The 60--85$\times$ slowdown that ruled
NUTS out in Python therefore \emph{does not transfer} to the
Julia port. The default Julia port still uses HMC (replicating the
framework-doc default), but NUTS becomes a viable future-work
candidate under §\ref{sec:in-out-julia}'s ``what changes when the
Python+JAX architectural constraints are removed'' frame.

\section{The port plan --- file-by-file mapping}
\label{sec:port-plan}

The Julia package is named \texttt{SMC2FC.jl} and lives at
\path{python-smc2-filtering-control/julia/SMC2FC/} (sister directory
to the Python \path{smc2fc/}). The port proceeds bottom-up so that
lower layers are tested before upper layers depend on them.

\subsection{Strategic principle: optimise for Julia, not a Python
re-skin}
\label{sec:port-principle}

The most common mistake an agent will make on this port is
translating Python line-by-line. The result is a wooden,
Python-shaped Julia package that realises none of Julia's actual
advantages --- and is therefore not worth the porting effort at
all. The point of doing this work is to land in a codebase that
is \emph{idiomatic Julia}: shorter, faster, more type-safe, and
better matched to the LEAN4 spec language than the Python
original. \textbf{A ``Python in Julia syntax'' port fails the
charter even if it type-checks and passes the differential test},
because it does not realise any of the §\ref{sec:why-julia}
benefits.

The translation discipline is therefore to port the
\emph{algorithm} and let Julia's native features (loops, multiple
dispatch, broadcasting, \texttt{AbstractArray} dispatch) express it
the Julia way:

\paragraph{Replace JAX control flow primitives with native Julia
loops.} \texttt{jax.lax.scan}, \texttt{jax.lax.fori\_loop},
\texttt{jax.lax.while\_loop} exist in JAX because Python's
interpreter overhead makes a plain \texttt{for} loop unusably slow.
\textbf{Julia has no such overhead.} A standard
\texttt{for k in 1:n} loop in Julia is fast under both ordinary
execution and \texttt{Enzyme.jl} / \texttt{Zygote.jl} autodiff. Do
not port \texttt{lax.scan} to a \texttt{Tullio.jl} reduction or to
a \texttt{KernelAbstractions.jl} kernel out of habit --- port it to
a \texttt{for} loop. The control-flow scaffolding that bloats the
JAX source files (the $\sim 5{,}600 \to \sim 1{,}200$ LOC reduction
from §\ref{sec:scaffolding} comes mostly from this) collapses into
ordinary Julia.

\paragraph{Replace conditional logic with multiple dispatch.}
Where Python branches on flags (e.g.\ the \texttt{bridge\_type}
switch in \path{core/sf_bridge.py}, the soft-vs-hard
chance-constrained switch in \path{control_v5.py}), Julia
dispatches on a small set of marker types:
\begin{lstlisting}
struct GaussianBridge end
struct SchrodingerFollmerBridge end

bridge_init(::GaussianBridge, prev_cloud, new_cloud) = ...
bridge_init(::SchrodingerFollmerBridge, prev_cloud, new_cloud) = ...
\end{lstlisting}
The conditional logic disappears into method definitions; each
behaviour lives in its own method, the dispatcher chooses at
compile time, and the agent cannot accidentally reach the wrong
branch.

\paragraph{Flatten nesting.} Python's \texttt{vmap}-over-particles
+ \texttt{lax.while\_loop}-over-tempering +
\texttt{lax.scan}-over-bins nests three primitives because each
layer needs its own JAX-recognised iterator. In Julia: nested
\texttt{for} loops, end of story. Broadcasting
(\texttt{f.(particles)}) handles the \texttt{vmap} layer; ordinary
loops handle the rest.

\paragraph{Use \texttt{AbstractArray\{T,N\}} in signatures, not
concrete \texttt{Array} or \texttt{CuArray}.} This is the
GPU-portability discipline from §\ref{sec:scaffolding} stated as a
translation rule: write the function once against the abstract
type, run it on whichever backend the caller supplied. The agent
should never need to write two versions of any function.

\paragraph{Zero-allocation inner loop --- opt-in mutability where it
counts.} JAX is pure-functional and relies on XLA to fuse and
eliminate the intermediate arrays each step would otherwise
allocate. Julia has no such fuser. A naive \texttt{for} loop that
returns a fresh \texttt{x\_new = f(x)} array on every iteration
\emph{will run substantially slower than the equivalent JAX
\texttt{lax.scan}}. The discipline is therefore: keep the
\emph{public API} functional (the \texttt{!} convention from
§\ref{sec:why-julia}.2 still holds at module boundaries), but in
the \emph{inner hot loops}, preallocate the working buffers once
and use in-place mutating operations:
\begin{lstlisting}
# Hot SDE step: NO per-iteration allocation.
function step!(x::AbstractVector, dt::Float64, theta::DynParams)
    x .= x .+ dt .* drift(x, theta)
    return x  # x has been mutated in place
end
\end{lstlisting}
A reasonable rule of thumb: every Phase 2--5 inner loop should be
benchmarked with \texttt{BenchmarkTools.@btime} and report
\texttt{0 allocations} per iteration (or a single allocation that is
amortised). Anything else is a missed optimisation, not a
shippable result.

\paragraph{Explicit CPU multithreading via
\texttt{Threads.@threads} or \texttt{Polyester.@batch}.} A subtlety
the previous bullets do not capture: JAX's \texttt{vmap} parallelises
the per-particle work on multi-core CPUs implicitly. A plain Julia
\texttt{for p in particles ... end} is sequential. To match (or beat)
the JAX-on-CPU baseline, outer particle loops have to be wrapped in a
threading macro:
\begin{lstlisting}
# Outer SMC loop, CPU multithreaded.
Threads.@threads for m in 1:n_smc
    log_lik[m] = inner_pf_loglikelihood(theta_dyn[m], data)
end

# Or, lower-overhead, for tight loops:
Polyester.@batch for m in 1:n_smc
    log_lik[m] = inner_pf_loglikelihood(theta_dyn[m], data)
end
\end{lstlisting}
On the GPU backend, \texttt{KernelAbstractions.@kernel} replaces
both. The agent should not ship any outer-particle loop in serial
form on the CPU backend.

\paragraph{Heavy broadcasting can hide allocations --- use
\texttt{@.} or explicit loops in hot paths.} The earlier
``broadcasting suffices'' note is true for simple operations, but
chained broadcasting on custom structs can quietly create
intermediate temporary arrays unless every operation is fused into a
single \texttt{@.} block:
\begin{lstlisting}
# Slow: each `.` may allocate a temporary.
x_new = a .+ b .* sin.(c) .- d

# Fast: single fused broadcast, no temporaries.
x_new = @. a + b * sin(c) - d

# Fastest in tight inner loops: explicit loop with @inbounds @simd.
@inbounds @simd for i in eachindex(x_new)
    x_new[i] = a[i] + b[i] * sin(c[i]) - d[i]
end
\end{lstlisting}
For complex inner-loop particle state transitions, the explicit
\texttt{@inbounds @simd} form (CPU) or a \texttt{KernelAbstractions}
kernel (GPU) is safer and usually faster than chained broadcasting.

\paragraph{Type stability is the \#1 Julia performance killer ---
\texttt{JET.jl} is the lint that catches it.} A function whose
return type the compiler cannot infer at compile time falls back to
runtime dynamic dispatch on every call. Inside a hot SMC$^2$ inner
loop, a single type-unstable line can cost 10--100$\times$ in
performance and is invisible at the source level. \texttt{JET.jl}'s
\texttt{report\_package} catches this. The audit checklist
(§\ref{sec:checklist-julia}) requires it green; if it flags a
\texttt{Union\{Float64, Nothing\}} return in a hot path, the agent
must refactor (e.g.\ replace \texttt{Nothing} sentinels with
\texttt{NaN}, narrow the input types, or hoist the conditional out
of the loop) before shipping.

These principles apply pervasively across the phases that follow.
The phase descriptions list \emph{what} to port; this subsection
is \emph{how} to port it.

\subsection{Phase 1 --- Foundations (\textasciitilde 1 hour)}

\begin{enumerate}
\item \texttt{Project.toml} + \texttt{Manifest.toml} with the
  dependency graph from §\ref{sec:scaffolding} (including
  \texttt{CUDA.jl} and \texttt{KernelAbstractions.jl}).
\item \texttt{src/Types.jl} --- the parametric types that encode
  the dimensional distinctions Python erases. Concretely:
  \begin{itemize}
  \item \texttt{State\{6\} = SVector\{6, Float64\}} (from
    \texttt{StaticArrays.jl}) --- stack-allocated, fixed-size,
    zero-overhead.
  \item \texttt{DynParams = ComponentVector\{Float64\}} (from
    \texttt{ComponentArrays.jl}) --- dot-access syntax
    (\texttt{theta.k\_FB}) on a flat \texttt{AbstractVector} that
    \texttt{AdvancedHMC.jl} and \texttt{Enzyme.jl} accept directly.
  \item \texttt{ParticleCloud = StructArray\{Particle\}} (from
    \texttt{StructArrays.jl}) --- struct-of-arrays layout for
    coalesced GPU access.
  \item \texttt{GPUFilterState\{T\}} and \texttt{CPUParameterCloud\{T\}}
    --- the hybrid-execution type boundary from
    §\ref{sec:hybrid}. The former wraps \texttt{CuArray} fields
    (state particles, weights); the latter is plain
    \texttt{Matrix\{T\}} / \texttt{Vector\{T\}}. Multiple
    dispatch chooses the right kernel based on which type is
    passed; the agent does not write two versions of any
    function.
  \item Marker types for dispatch: \texttt{struct GaussianBridge
    end}, \texttt{struct SchrodingerFollmerBridge end},
    \texttt{struct SoftSurrogate end}, \texttt{struct
    HardIndicator end}.
  \end{itemize}
  All variable-length array fields use the abstract
  \texttt{AbstractArray\{T,N\}} where they could be either CPU
  (\texttt{Array}) or GPU (\texttt{CuArray}); fields that the
  hybrid contract pins to a specific device use the concrete
  type so the dispatcher can route correctly. This composite
  type-system is the LEAN4-charter-spirit applied to the
  framework.
\item \texttt{src/Config.jl} --- translation of
  \path{core/config.py} \texttt{SMCConfig} as a \texttt{@kwdef}
  struct, plus a \texttt{backend::Symbol = :cpu} field
  (\texttt{:cpu} or \texttt{:cuda}) that selects the array
  constructor.
\item \texttt{src/Transforms.jl} --- wraps \texttt{Bijectors.jl}
  for the unconstrained-space mapping.
\end{enumerate}

\subsection{Phase 2 --- Filtering (\textasciitilde 2 hours, GPU-resident)}

Phase 2 is the inner SMC$^2$ loop and per the hybrid strategy of
§\ref{sec:hybrid} it runs \emph{on the GPU}: state particles are
\texttt{CuArray}-resident and the kernels dispatch to CUDA via
\texttt{StochasticDiffEq.jl} or \texttt{KernelAbstractions.@kernel}.

\begin{enumerate}\setcounter{enumi}{4}
\item \texttt{src/Filtering/Kernels.jl} --- Liu--West shrinkage
  moment match, Gaussian smoothing kernels, ESS computation. CUDA-
  dispatched via \texttt{AbstractArray\{T,N\}}-generic code.
\item \texttt{src/Filtering/Bootstrap.jl} --- the inner bootstrap
  filter, producing $\widehat{L}_N(\theta_{\rm dyn})$. The output
  is a single scalar (or a small vector if computing a gradient)
  that crosses back to the host CPU via PCIe; particle states
  stay on the GPU.
\item \texttt{src/Filtering/OT.jl} --- wraps
  \texttt{OptimalTransport.jl} for the OT rescue; implements the
  sigmoid blend with the systematic-resample output.
\item Tests for Phase 2: differential against Python on a battery
  of random SDE+observation pairs. Tolerance $10^{-6}$ on per-step
  weights, $10^{-4}$ on integrated likelihood. \textbf{Tests run
  twice: once with \texttt{CPU()} arrays, once with
  \texttt{CUDA.CuArray}; both must pass within the same
  tolerance.}
\end{enumerate}

\subsection{Phase 3 --- Outer SMC$^2$ (\textasciitilde 1.5 hours, CPU-resident)}

Phase 3 is the outer SMC$^2$ loop and per §\ref{sec:hybrid} it
runs \emph{on the CPU}: parameter cloud, log-likelihood vector,
mass matrix, and HMC trajectories are all
\texttt{Array}-resident. The likelihood callable provided to
\texttt{AdvancedHMC.jl} performs a host-to-device transfer of
the 37-D parameter vector and a device-to-host transfer of the
scalar log-likelihood per gradient evaluation; both are
microsecond-scale across PCIe Gen 5.

\begin{enumerate}\setcounter{enumi}{8}
\item \texttt{src/SMC2/Tempering.jl} --- adaptive
  $\delta\lambda$ via bisection on CPU. Direct port of
  \texttt{\_solve\_delta\_for\_ess}; the JAX
  \texttt{lax.while\_loop} of the Python becomes a plain Julia
  \texttt{while} loop per §\ref{sec:port-principle}.
\item \texttt{src/SMC2/HMC.jl} --- wraps \texttt{AdvancedHMC.jl}
  for the rejuvenation moves on the CPU's branch-predictor-rich
  cores; re-estimates the diagonal mass matrix per tempering
  level. The likelihood closure dispatches into Phase 2's GPU
  filter for evaluation.
\item \texttt{src/SMC2/TemperedSMC.jl} --- the outer loop, calling
  Phase 2 for likelihood, Phase 3.1--3.2 for tempering and HMC.
\item \texttt{src/SMC2/Bridge.jl} --- Gaussian + BW-geodesic
  warm-starts, dispatched on marker types per
  §\ref{sec:port-principle}.
\item Tests for Phase 3: differential against Python on the same
  benchmark \path{tests/test_obs_consistency_v5.py} already uses;
  both backends.
\end{enumerate}

\subsection{Phase 4 --- Control (\textasciitilde 1 hour)}

\begin{enumerate}\setcounter{enumi}{13}
\item \texttt{src/Control/Spec.jl} --- \texttt{ControlSpec}
  struct (cost callable, decoder, horizon, initial state).
\item \texttt{src/Control/Calibration.jl} --- $\beta_{\max}$
  auto-calibration.
\item \texttt{src/Control/RBFSchedule.jl} --- RBF schedule
  decoder.
\item \texttt{src/Control/TemperedSMC.jl} --- re-uses Phase 3.3
  with the cost-as-likelihood substitution.
\item Tests for Phase 4: differential against Python on the
  chance-constrained controller benchmark; both backends.
\end{enumerate}

\subsection{Phase 5 --- Plant + Simulator (\textasciitilde 0.5
hour)}

\begin{enumerate}\setcounter{enumi}{18}
\item \texttt{src/Simulator/SDEModel.jl} --- wraps
  \texttt{StochasticDiffEq.jl}.
\item \texttt{src/Simulator/Observations.jl} --- per-channel
  observation sampling.
\end{enumerate}

\subsection{Phase 6 --- End-to-end + GPU benchmark
(\textasciitilde 1.5 hours)}

\begin{enumerate}\setcounter{enumi}{20}
\item Full-pipeline test: run a complete inference + control loop
  on the FSA-v5 model (consuming the LEAN4 reference from Part I),
  compare end-to-end against the Python pipeline on the same
  inputs. Run on \textbf{both backends}
  (\texttt{backend = :cpu} and \texttt{backend = :cuda}); confirm
  numerical agreement to within the same tolerances as Phase 2.
\item \texttt{JET.jl} static analysis pass; fix any flagged type
  uncertainties.
\item Benchmark per-window cost on CPU and GPU against the
  JAX-native Python baseline of $\sim 1$ second/window. Record
  the numbers in the package README; this is the artefact that
  justifies the port to anyone reading it later.
\end{enumerate}

\section{Scope and timing (framework side)}
\label{sec:scope-julia}

\textbf{This is translation between two scientific computing
languages.} The user has explicitly rejected inflated estimates.
With the existing Julia scaffolding from §\ref{sec:scaffolding}
doing the heavy lifting, the port is:

\begin{center}\small
\begin{tabular}{ll}
\toprule
Phase & Estimate \\
\midrule
Phase 1 --- Foundations                        & 1 hour \\
Phase 2 --- Filtering (CPU + GPU)              & 2 hours \\
Phase 3 --- Outer SMC$^2$ (CPU + GPU)          & 1.5 hours \\
Phase 4 --- Control (CPU + GPU)                & 1 hour \\
Phase 5 --- Plant + Simulator (CPU + GPU)      & 0.5 hour \\
Phase 6 --- End-to-end + GPU benchmark         & 1.5 hours \\
\midrule
\textbf{Total}                                 & \textbf{$\sim 7.5$ hours of focused work} \\
\bottomrule
\end{tabular}
\end{center}

\textbf{The CPU/GPU dual-target adds essentially nothing} to the
per-phase budgets above because Julia's generic dispatch on
\texttt{AbstractArray\{T,N\}} means the same code runs on both
\texttt{Array} (CPU) and \texttt{CuArray} (GPU); the library
packages (\texttt{AdvancedHMC.jl}, \texttt{StochasticDiffEq.jl},
\texttt{OptimalTransport.jl}, \texttt{Distributions.jl},
\texttt{Bijectors.jl}) already have GPU support upstream. The
marginal GPU cost is in (a) running the existing test suite a
second time on \texttt{CuArray}, and (b) the Phase 6 benchmark
--- both already accounted for in the table.

For a developer competent in both languages, this is a single
focused day's work. Most of it is wiring library calls together
and running differential tests; the algorithmic content is not
re-derived but transcribed. The resulting Julia is markedly
shorter ($\sim 1{,}200$ LOC vs $\sim 5{,}600$ LOC) because the
boilerplate that Python+JAX requires (closures, \texttt{Partial}
wrappers, \texttt{lax.while\_loop} scaffolding, manual
\texttt{jit} placement) collapses in Julia.

\textbf{The 7.5-hour estimate assumes an idiomatic, performance-
optimised Julia port, not a line-by-line Python re-skin.} A wooden
translation that ports each \texttt{lax.scan} to a \texttt{Tullio}
reduction and keeps Python-shaped control flow would technically
compile in less time, but it would not realise any of the
§\ref{sec:why-julia} benefits and would fail the
§\ref{sec:port-principle} discipline. The estimate above is for
code written the Julia way: plain \texttt{for} loops, multiple
dispatch on marker types, \texttt{AbstractArray\{T,N\}}
signatures, \texttt{StaticArrays.SVector} for fixed-size objects,
\texttt{StructArrays} where the natural memory layout is
struct-of-arrays, \texttt{Threads.@threads} or
\texttt{Polyester.@batch} on outer particle loops, and
\texttt{JET.jl}-clean type stability throughout. The
performance-discipline patterns of §\ref{sec:port-principle}
shape \emph{how} the budgeted hours are spent; they do not add
hours, because the right Julia idioms are typically shorter and
faster than their Python+JAX equivalents (the
$\sim 5{,}600 \to \sim 1{,}200$ LOC reduction is exactly that
shrinkage). If the agent finds itself porting an algorithm faster
than the table predicts, the result should still be checked
against §\ref{sec:port-principle} --- speed of writing is not the
objective; an idiomatic, type-safe, GPU-portable, zero-allocation-
inner-loop Julia codebase is.

The cost/value frame: a single 7.5-hour port amortises across
\emph{every future model the framework runs} (the JAX-native vs
BlackJAX backend split disappears; the per-window recompile
landmine disappears; the NUTS-inside-vmap cliff disappears; the
agent pattern-match risk on framework internals disappears). The
current trajectory --- paying repeatedly for agent rounds of
pattern-match-then-debug-Python+JAX-internals --- is what this
charter exists to stop.

\section{What is and isn't covered (framework side)}
\label{sec:in-out-julia}

\paragraph{In scope.} A Julia package \texttt{SMC2FC.jl}
providing filtering (bootstrap PF + OT rescue), outer SMC$^2$
(tempered + HMC + adaptive mass matrix + bridges),
control-as-inference, an SDE simulator, and the
unconstrained-space transforms. \textbf{Both CPU and GPU (CUDA)
backends are in scope from day one}, and the production
configuration is the \textbf{hybrid pipeline of
§\ref{sec:hybrid}}: GPU for the inner particle filter, CPU for
the outer SMC$^2$ rejuvenation. CPU-only and GPU-only
configurations are also tested as diagnostic baselines.
Differential tests against the Python implementation on the
existing benchmarks for all configurations.

\paragraph{Out of scope (Python stays in service).} The Python
\texttt{smc2fc/} package continues to exist during and after the
port. The Julia port does \emph{not} replace it; it sits
alongside as a type-checked alternative implementation. Existing
user code that imports the Python package is unaffected.

\paragraph{Out of scope for the default port (settled in framework
doc).} The hard-vs-soft chance-constrained finding (use soft
only) and the SF/BW-vs-Gaussian-bridge equivalence finding carry
over from Python; the Julia port replicates the soft-HMC +
Gaussian-bridge defaults.

\paragraph{Re-opened by the hybrid architecture (future work).}
The NUTS-as-HMC-replacement finding from framework-doc §6.4 is a
\emph{Python+JAX architectural} result: NUTS-inside-vmap-inside-
\texttt{lax.while\_loop} suffers warp divergence and dynamic-tree-
expansion costs that combine for a 60--85$\times$ slowdown. The
Julia hybrid pipeline of §\ref{sec:hybrid} runs NUTS on the CPU
where its branchy control flow is well-matched to the hardware,
which means \emph{the empirical reason the framework rejected NUTS
does not transfer to Julia}. The default Julia port still uses
HMC for parity with the Python implementation, but a NUTS-on-
hybrid benchmark is a candidate for the Phase 6 follow-on once the
HMC port is green.

\section{Audit checklist (framework-side agent)}
\label{sec:checklist-julia}

A future agent runs this at session start, before touching any
Julia file:

\begin{itemize}
\item[$\square$] Is there a Python reference for the function I
  am about to port?
\item[$\square$] Does the LEAN4-first charter's process (Part I)
  apply (i.e.\ am I porting model code that should consume a
  LEAN4 spec)? If yes, does the spec exist?
\item[$\square$] Have I declared parametric types for every
  input and output, capturing the dimensional information Python
  erases?
\item[$\square$] Have I left array fields as
  \texttt{AbstractArray\{T,N\}} so the same code accepts
  \texttt{Array} (CPU) and \texttt{CuArray} (GPU)?
\item[$\square$] Did I port the \emph{algorithm}, not the
  \emph{code}? Specifically: no \texttt{jax.lax.scan} /
  \texttt{lax.fori\_loop} / \texttt{lax.while\_loop} analogues
  survived the translation --- they became plain Julia
  \texttt{for} loops. (Per §\ref{sec:port-principle}.)
\item[$\square$] Did I express type-dispatched behaviours via
  multiple dispatch (one method per marker struct) rather than
  as flag-driven conditionals carried over from Python? (Per
  §\ref{sec:port-principle}.)
\item[$\square$] Did I use \texttt{StaticArrays.SVector} for
  fixed-size small objects (\texttt{State\{6\}}, parameter
  blocks, RBF coefficient vectors) rather than \texttt{Vector}?
  (Heap-allocated \texttt{Vector}s in the inner loop regress
  versus JAX.)
\item[$\square$] Did I use \texttt{StructArrays} or
  \texttt{ComponentArrays} where the natural memory layout is
  Struct-of-Arrays (particle clouds, parameter blocks)?
\item[$\square$] Are inner loops zero-allocation?
  Verify with \texttt{BenchmarkTools.@btime}; an inner SDE step,
  Liu--West shrinkage step, or HMC leapfrog step should report
  \texttt{0 allocations} per iteration.
\item[$\square$] Are outer particle loops parallelised on the
  CPU backend (\texttt{Threads.@threads} or
  \texttt{Polyester.@batch})? On the GPU backend, replaced by a
  \texttt{KernelAbstractions} kernel?
\item[$\square$] Does the function I just ported respect the
  hybrid contract from §\ref{sec:hybrid}? Specifically: does
  inner-loop SDE/PF work stay on the GPU (\texttt{CuArray}
  fields), does outer-loop HMC/tempering work stay on the
  CPU (\texttt{Array} fields), and is the only host--device
  transfer the scalar log-likelihood (or the small parameter
  vector going the other way)?
\item[$\square$] If I wrote a likelihood callable that
  \texttt{AdvancedHMC.jl} consumes, does it explicitly perform
  the host-to-device transfer on the parameter vector and the
  device-to-host transfer on the scalar log-likelihood, with
  nothing else crossing the bus per evaluation?
\item[$\square$] Does
  \texttt{julia --project=. -e 'using Pkg; Pkg.test()'} pass on
  \textbf{both} the CPU and the CUDA test paths?
\item[$\square$] Does
  \texttt{julia --project=. -e 'using JET; report\_package(SMC2FC)'}
  report no type uncertainties \emph{or} type instabilities?
  (Type instability is the \#1 Julia performance killer; flag
  any \texttt{Union\{...\}} return in a hot inner loop.)
\item[$\square$] Does the differential test against the Python
  implementation pass within the tolerance threshold for the
  function I just ported, on both backends?
\item[$\square$] Have I updated the §\ref{sec:port-plan} port-mapping
  table to mark this file as ported?
\end{itemize}

If any box is unchecked, the agent does not proceed.

\section{Closing of Part II}

The point of Part II is accountability on the framework side.
Python+JAX is a perfect storm for the kind of mathematical work
\texttt{smc2fc} does: untyped arrays, invisible JIT contracts,
closures over data, nesting cliffs. LLM agents pattern-matching
through that surface burn the user's time and tokens hunting the
consequences. Julia narrows the surface by structural means --- a
single command (\texttt{julia --project=. -e 'using JET;
report\_package(SMC2FC)'}) catches the kind of error that takes a
human reviewer hours to find in the JAX/Python stack.

Combined with Part I, this charter establishes a single typed
pipeline:

\begin{center}
\fbox{\parbox{0.9\textwidth}{\centering
\textbf{LaTeX equations} $\to$
\textbf{LEAN4 reference} (model side, Part I) $\to$
\textbf{Julia framework} (\texttt{smc2fc} side, Part II)\\[0.3em]
\emph{with Python remaining in service as the differentially-tested
fast path on both sides.}
}}
\end{center}

LEAN4 is never wrong about types. Julia narrows the dispatch surface
to the point that pattern-matching does not silently produce wrong
results. Python is fast and well-tested but pattern-match-friendly
to a fault; it stays in service \emph{as the fast path that gets
checked against the typed reference}, not as the source of truth.

The total commitment in human hours --- $\sim 8$ for Part I,
$\sim 7.5$ for Part II, single-day-each scale --- is the price
of replacing repeated multi-day rounds of agent pattern-match-then-
debug with a one-shot translation that holds for every future
model and every future agent.

\end{document}
